\chapter{Additive Functions and the Kubilius Model}\label{ch:additive}
\begin{center}\itshape
Pick an integer at random. Is it abundant? How close is $\varphi(n)$ to $n$?\\
The primes dividing $n$ behave like independent coins,\\
and an additive function is the sum of what they show.
\end{center}

In 1928 Isaac Schoenberg, a young Romanian mathematician then working in Germany, proved a curious fact about Euler's function: the proportion of integers $n\le x$ with $\frac{\varphi(n)}n\le z$ tends, for every $z$, to a limit $F(z)$, and $F$ is continuous. Five years later Harold Davenport proved the same kind of statement for $\frac{\sigma(n)}n$, the sum of the divisors of $n$ divided by $n$. As a consequence the \emph{abundant numbers}, those with $\sigma(n)\ge2n$, have a natural density, a question that went back to the Greeks' fascination with perfect numbers. In the same year Felix Behrend showed that the density lies between $0.241$ and $0.314$. It is now known to be $0.24761\ldots$, which means that roughly one integer in four is abundant.

These were hard-won theorems, each proved with its own machinery. In 1939 Paul Erd\H{o}s and Aurel Wintner, in a nine-page paper in the \emph{American Journal of Mathematics}, explained them all at once. Write $\log\frac{\varphi(n)}n=\sum_{p\mid n}\log(1-\frac1p)$. This is an \emph{additive} function: a sum over the primes dividing $n$. Erd\H{o}s and Wintner gave a criterion for an additive function to have a limiting distribution, and the criterion is the three-series theorem that Kolmogorov had proved in 1928 for sums of independent random variables. The title of their paper says why: \emph{Additive arithmetical functions and statistical independence}.

The independence is a theorem and not just a heuristic. Twenty years later Jonas Kubilius in Vilnius made it precise: the exponents of the small primes in a random integer $n\le x$ are, up to a tiny error in total variation, independent geometric random variables. Kubilius became rector of Vilnius University in 1958 and led it for more than thirty years, until 1990, while building the Lithuanian school of probability and number theory. His monograph, translated by the American Mathematical Society in 1964 as \emph{Probabilistic Methods in the Theory of Numbers}, gave the subject its name.

This chapter builds the model and draws the consequences: the Tur\'an--Kubilius inequality for every additive function, the Kubilius model in total variation, the Erd\H{o}s--Wintner theorem, and two theorems of probability, of L\'evy and of Jessen and Wintner, which decide whether the limit law is continuous. Euler's function and the abundant numbers come out as corollaries.

\begin{reminder}[from Tomes~3, 5, 10 and~13]
\emph{Tur\'an--Kubilius} (Tome~5, Theorem~9.1). There is an absolute $C_0$ such that for every real function $f$ on the primes and $x\ge3$, the strongly additive $f(n)=\sum_{p\mid n}f(p)$ satisfies $\sum_{n\le x}(f(n)-A_0(x))^2\le C_0xB_0(x)^2$, where $A_0(x)=\sum_{p\le x}\frac{f(p)}p$ and $B_0(x)^2=\sum_{p\le x}\frac{f(p)^2}p$.
\emph{Rough numbers} (Tome~13, Corollary~3.5). For $x\ge4$ and $2\le z\le\sqrt x$, the number $\Phi(x,z)$ of $n\le x$ with no prime factor $\le z$ is $x\prod_{p\le z}(1-\frac1p)(1+O(e^{-u/2}))+O(\sqrt x\log x)$, with $u=\frac{\log x}{\log z}$.
\emph{Three-series theorem} (Tome~10, Theorem~2.5). For independent $X_p$ and $c>0$, $\sum X_p$ converges almost surely if and only if $\sum P(|X_p|>c)$, $\sum\mathbb E[X_p\mathbf 1_{|X_p|\le c}]$ and $\sum\operatorname{Var}[X_p\mathbf 1_{|X_p|\le c}]$ converge.
\emph{Kolmogorov's zero--one law} (Tome~10, Theorem~1.4) and the \emph{Borel--Cantelli lemmas} (Tome~3, Chapter~4).
\emph{Mertens} (Tome~5, Theorems~1.3--1.6). $\sum_{p\le y}\frac{\log p}p=\log y+O(1)$, $\sum_{p\le y}\frac1p=\log\log y+O(1)$ and $\prod_{p\le y}(1-\frac1p)\asymp\frac1{\log y}$.
\end{reminder}

\section{Additive functions}

\begin{definition}\label{def:additive}
A function $f:\N_{\ge1}\to\C$ is \emph{additive} if $f(mn)=f(m)+f(n)$ whenever $\gcd(m,n)=1$, so that $f(n)=\sum_{p^k\parallel n}f(p^k)$, where $p^k\parallel n$ means that $p^k$ divides $n$ and $p^{k+1}$ does not. It is \emph{strongly additive} if moreover $f(p^k)=f(p)$ for all $k\ge1$, and \emph{completely additive} if $f(mn)=f(m)+f(n)$ for all $m,n$.
\end{definition}

The number $\omega(n)$ of distinct prime factors is strongly additive, $\Omega(n)$ (with multiplicity) and $\log n$ are completely additive, and $\log\frac{\varphi(n)}n=\sum_{p\mid n}\log(1-\frac1p)$ is strongly additive. If $g$ is multiplicative and positive, $\log g$ is additive; so $\log\frac{\sigma(n)}n$ is additive, with $\log\frac{\sigma(p^k)}{p^k}=\log\frac{1-p^{-k-1}}{1-p^{-1}}$.

The object of study is the \emph{distribution} of $f$: for $n$ chosen uniformly in $\{1,\dots,\lfloor x\rfloor\}$, the law of the random variable $f(n)$. We write $P_x$ and $\mathbb E_x$ for probabilities and expectations under this uniform choice, so that $P_x(f(n)\le z)=\frac1{\lfloor x\rfloor}\#\{n\le x:f(n)\le z\}$. We say that $f$ has a \emph{limiting distribution} if $f(n)$ converges in law as $x\to\infty$, that is, if $P_x(f(n)\le z)$ tends to $F(z)$ at every continuity point $z$ of some distribution function $F$.

\begin{think}
Which of $\omega(n)$, $\log\frac{\varphi(n)}n$ and $\log\frac{\sigma(n)}n$ would you expect to have a limiting distribution without any renormalisation? Think of the size of the contribution of each prime, and of how often that prime divides $n$.
\end{think}

The heuristic is simple. A prime $p$ divides a random integer with probability $\frac1p$, and does so ``independently'' of the other primes. For $\omega$ each prime contributes $1$ with probability $\frac1p$, and $\sum\frac1p$ diverges: $\omega(n)$ grows like $\log\log n$ and must be centred and scaled, which is the Erd\H{o}s--Kac theorem of Tome~5 and of Chapter~3. For $\log\frac{\varphi(n)}n$ the prime $p$ contributes about $-\frac1p$ with probability $\frac1p$, and $\sum\frac1{p^2}$ converges: the sum should converge without any normalisation. The rest of the chapter makes this precise.

\section{The Tur\'an--Kubilius inequality}

Tome~5 proved the inequality for strongly additive functions. Every additive function differs from a strongly additive one only at the prime powers $p^k$ with $k\ge2$, which are rare, and that is the whole proof of the general case.

\begin{theorem}[Tur\'an--Kubilius]\label{thm:tkgeneral}
There is an absolute constant $C_1$ such that for every additive $f:\N\to\C$ and every $x\ge3$,
\[
\sum_{n\le x}|f(n)-A(x)|^2\le C_1xB(x)^2,\qquad A(x)=\sum_{p^k\le x}\frac{f(p^k)}{p^k}\Bigl(1-\frac1p\Bigr),\qq B(x)^2=\sum_{p^k\le x}\frac{|f(p^k)|^2}{p^k}.
\]
One can take $C_1=3C_0+36$.
\end{theorem}

The centring $A(x)$ is the expectation of the \emph{model}: if $G_p$ is a geometric variable with $P(G_p=k)=(1-\frac1p)p^{-k}$, then $\mathbb E f(p^{G_p})=\sum_{k\ge1}f(p^k)(1-\frac1p)p^{-k}$, with $f(1)=0$.

\begin{proof}
Since $|w|^2=(\operatorname{Re}w)^2+(\operatorname{Im}w)^2$ and the real and imaginary parts of $f$ are additive with $A$ and $B^2$ splitting accordingly ($B^2=B_{\operatorname{Re}}^2+B_{\operatorname{Im}}^2$), it suffices to treat real $f$ with the constant $3C_0+36$.

For $n\le x$ write $f(n)=g(n)+h(n)$ with
\[
g(n)=\sum_{p\mid n}f(p),\qquad h(n)=\sum_{\substack{p^k\parallel n\\k\ge2}}c(p^k),\qquad c(p^k)=f(p^k)-f(p).
\]
The function $g$ is strongly additive, and Tome~5, Theorem~9.1, gives $\sum_{n\le x}(g(n)-A_0)^2\le C_0xB_0^2\le C_0xB^2$, with $A_0=\sum_{p\le x}\frac{f(p)}p$.

\emph{The function $h$.} Expanding the square, the diagonal terms are $\sum_{p^k}c(p^k)^2\#\{n\le x:p^k\parallel n\}\le x\sum_{p^k}\frac{c(p^k)^2}{p^k}$. In the off-diagonal terms $p^k\parallel n$ and $q^l\parallel n$ with $p^k\ne q^l$ force $p\ne q$, so $p^kq^l\mid n$ and there are at most $\frac x{p^kq^l}$ such $n$. Hence, with all sums over $p^k\le x$ with $k\ge2$, and by Cauchy--Schwarz,
\[
\sum_{n\le x}h(n)^2\le x\sum\frac{c(p^k)^2}{p^k}+x\Bigl(\sum\frac{|c(p^k)|}{p^k}\Bigr)^2\le x\,H\Bigl(1+\sum_p\sum_{k\ge2}\frac1{p^k}\Bigr)\le2xH,
\]
where $H=\sum\frac{c(p^k)^2}{p^k}$, because $\sum_p\sum_{k\ge2}p^{-k}=\sum_p\frac1{p(p-1)}<\sum_{m\ge2}\frac1{m(m-1)}=1$. Now $c(p^k)^2\le2f(p^k)^2+2f(p)^2$, and $\sum_{k\ge2}p^{-k}=\frac1{p(p-1)}\le\frac1p$, so
\[
H\le2\sum_{\substack{p^k\le x\\k\ge2}}\frac{f(p^k)^2}{p^k}+2\sum_{p\le x}\frac{f(p)^2}p\le4B^2,\qquad\text{and}\qquad\sum_{n\le x}h(n)^2\le8xB^2.
\]

\emph{The centring.} Since $\frac{f(p)}p-\frac{f(p)}p(1-\frac1p)=\frac{f(p)}{p^2}$,
\[
A_0-A=\sum_{p\le x}\frac{f(p)}{p^2}-\sum_{\substack{p^k\le x\\k\ge2}}\frac{f(p^k)}{p^k}\Bigl(1-\frac1p\Bigr),
\]
and by Cauchy--Schwarz each sum is at most $B$ in absolute value: $|\sum_p\frac{f(p)}{p^2}|\le(\sum_p\frac{f(p)^2}p)^{1/2}(\sum_p\frac1{p^3})^{1/2}\le B$ and $|\sum\frac{f(p^k)}{p^k}|\le B(\sum_p\sum_{k\ge2}p^{-k})^{1/2}\le B$. So $|A_0-A|\le2B$.

Finally $f(n)-A=(g(n)-A_0)+h(n)+(A_0-A)$ and $(a+b+c)^2\le3(a^2+b^2+c^2)$ give $\sum_{n\le x}(f(n)-A)^2\le3(C_0+8+4)xB^2$.
\end{proof}

In the language of probability: $\operatorname{Var}_xf(n)\le\mathbb E_x|f(n)-A(x)|^2\le C_1B(x)^2(1+O(\frac1x))$, and $B(x)^2$ is, up to bounded factors, the variance of the model $\sum_{p\le x}f(p^{G_p})$. The dependence between the primes costs at most a constant factor. Chebyshev's inequality turns this into a law of large numbers: if $B(x)\to\infty$, then $\frac{f(n)-A(x)}{B(x)\,\psi(x)}\to0$ in probability for every $\psi(x)\to\infty$. For $\omega$ this is the theorem of Hardy and Ramanujan (Tome~5, Theorem~9.2).

\section{The Kubilius model}

The inequality compares only two moments. Kubilius's insight was that much more is true: for the primes up to a small power of $x$, the whole joint law of the exponents is close to the independent model, in the strongest sense available.

For $y\ge2$ and $n\ge1$ let $v_y(n)=(\nu_p(n))_{p\le y}$ be the vector of exponents of the primes $p\le y$ in $n$. Let $(G_p)_p$ be independent with $P(G_p=k)=(1-\frac1p)p^{-k}$ for $k\ge0$, and $G^{(y)}=(G_p)_{p\le y}$. Recall that the \emph{total variation distance} between two laws $\mu,\nu$ on a countable set is $d_{TV}(\mu,\nu)=\sup_E|\mu(E)-\nu(E)|=\frac12\sum_a|\mu(a)-\nu(a)|$.

\begin{theorem}[Kubilius]\label{thm:kubilius}
There is an absolute constant $C_2$ such that for $x\ge16$ and $2\le y\le x^{1/4}$, with $u=\frac{\log x}{\log y}$,
\[
d_{TV}\bigl(\operatorname{law}_x v_y(n),\ \operatorname{law}G^{(y)}\bigr)\le C_2\bigl(e^{-u/4}+x^{-1/4}\log x\bigr).
\]
\end{theorem}

\begin{proof}
A vector $\mathbf k=(k_p)_{p\le y}$ with finitely many non-zero entries corresponds to a $y$-\emph{smooth} integer $m=\prod_pp^{k_p}$, one all of whose prime factors are $\le y$. Write $V=\prod_{p\le y}(1-\frac1p)$, so $P(G^{(y)}=\mathbf k)=\frac Vm$. On the other side, $v_y(n)=\mathbf k$ means $n=mn'$ with $n'$ free of prime factors $\le y$, so
\[
P_x(v_y(n)=\mathbf k)=\frac{\Phi(x/m,y)}{\lfloor x\rfloor}.
\]
\emph{Small $m$.} Let $m\le\sqrt x$. Then $X=\frac xm\ge\sqrt x\ge y^2$, and Tome~13, Corollary~3.5, applies with $u_X=\frac{\log X}{\log y}\ge\frac u2$:
\[
\Bigl|\frac{\Phi(x/m,y)}x-\frac Vm\Bigr|\ll\frac Vme^{-u/4}+\frac{\sqrt{x/m}\log x}x.
\]
Summing over $y$-smooth $m\le\sqrt x$, and using $\sum_{m\ y\text{-smooth}}\frac1m=\prod_{p\le y}(1-\frac1p)^{-1}=\frac1V$ and $\sum_{m\le\sqrt x}m^{-1/2}\le2x^{1/4}$,
\[
\Sigma_1:=\sum_{m\le\sqrt x}\Bigl|\frac{\Phi(x/m,y)}x-\frac Vm\Bigr|\ll e^{-u/4}+x^{-1/4}\log x.
\]
\emph{Large $m$, model side (Rankin's trick).} Let $\delta=\frac1{\log y}$. For $m>\sqrt x$, $1\le(\frac m{\sqrt x})^\delta$, so
\[
\Sigma_2:=\sum_{\substack{m>\sqrt x\\y\text{-smooth}}}\frac Vm\le Vx^{-\delta/2}\sum_{y\text{-smooth}}\frac1{m^{1-\delta}}=Ve^{-u/2}\prod_{p\le y}\Bigl(1-\frac1{p^{1-\delta}}\Bigr)^{-1}.
\]
For $p\le y$, $p^\delta=e^{\log p/\log y}\le1+(e-1)\frac{\log p}{\log y}$ by convexity, so $\sum_{p\le y}p^{\delta-1}\le\sum_{p\le y}\frac1p+\frac{e-1}{\log y}\sum_{p\le y}\frac{\log p}p=\log\log y+O(1)$, and $\sum_pp^{2\delta-2}\le\sum_pe^2p^{-2}=O(1)$. With $\log(1-t)^{-1}\le t+t^2$ for $0\le t\le\frac12$ (and the factor at $p=2$ bounded), the product is $\ll\log y$, and $V\ll\frac1{\log y}$. Hence $\Sigma_2\ll e^{-u/2}$.

\emph{Large $m$, integer side.} The probabilities $P_x(v_y(n)=\mathbf k)$ sum to $1$, and so do the $\frac Vm$. Since $\lfloor x\rfloor\le x$,
\[
\sum_{m>\sqrt x}\frac{\Phi(x/m,y)}{\lfloor x\rfloor}=1-\sum_{m\le\sqrt x}\frac{\Phi(x/m,y)}{\lfloor x\rfloor}\le1-\sum_{m\le\sqrt x}\frac{\Phi(x/m,y)}x\le1-\sum_{m\le\sqrt x}\frac Vm+\Sigma_1=\Sigma_2+\Sigma_1.
\]
\emph{Conclusion.} Since $\sum_{m\le\sqrt x}\Phi(x/m,y)\le x$, replacing $\frac1x$ by $\frac1{\lfloor x\rfloor}$ in $\Sigma_1$ costs at most $\frac x{\lfloor x\rfloor}-1\le\frac2x$. So
\[
2d_{TV}=\sum_{\mathbf k}\bigl|P_x(v_y(n)=\mathbf k)-P(G^{(y)}=\mathbf k)\bigr|\le\Bigl(\Sigma_1+\frac2x\Bigr)+(\Sigma_1+\Sigma_2)+\Sigma_2\ll e^{-u/4}+x^{-1/4}\log x.
\]
\end{proof}

\begin{think}
With $y=x^{1/u}$, the theorem is useful only when $u$ is large. What goes wrong for $y=\sqrt x$? (Count the $n\le x$ that are prime, and compare with the model's probability that $G_p=0$ for every $p\le\sqrt x$.)
\end{think}

For $y=\sqrt x$ the integer side has $P_x(v_y(n)=\mathbf0)=\frac{1+\pi(x)-\pi(\sqrt x)}{\lfloor x\rfloor}\sim\frac1{\log x}$, while the model gives $V\sim\frac{2e^{-\gamma}}{\log x}$: the two differ by the factor $\frac{e^\gamma}2\approx0.89$ already met in Tome~13, Figure~3.1. A large prime factor uses up most of the room below $x$, and that constraint is invisible to independent coins. The model is exact for the primes up to $x^{o(1)}$, and only for those.

Since $f(n)$ restricted to the small primes is a function of $v_y(n)$, the theorem transfers at once to every additive function. Let
\[
f_y(n)=\sum_{\substack{p^k\parallel n\\p\le y}}f(p^k),\qquad Z_y=\sum_{p\le y}f(p^{G_p})\quad(f(1)=0).
\]

\begin{corollary}\label{cor:kubiliusf}
For every additive $f$, $x\ge16$ and $2\le y\le x^{1/4}$, and every Borel set $E\subseteq\C$, $|P_x(f_y(n)\in E)-P(Z_y\in E)|\le C_2(e^{-u/4}+x^{-1/4}\log x)$.
\end{corollary}
\begin{proof}
$f_y(n)=\Lambda(v_y(n))$ and $Z_y=\Lambda(G^{(y)})$ for the same map $\Lambda(\mathbf k)=\sum_pf(p^{k_p})$, and total variation distance can only decrease under a map: $\{f_y(n)\in E\}=\{v_y(n)\in\Lambda^{-1}(E)\}$.
\end{proof}

\section{The Erd\H{o}s--Wintner theorem}

For a real additive $f$ and a prime $p$ write $f^*(p)=f(p)$ if $|f(p)|\le1$ and $f^*(p)=0$ otherwise.

\begin{theorem}[Erd\H{o}s and Wintner, 1939]\label{thm:erdoswintner}
Let $f$ be a real additive function such that the three series
\[
\sum_{|f(p)|>1}\frac1p,\qquad\sum_p\frac{f^*(p)}p,\qquad\sum_p\frac{f^*(p)^2}p
\]
converge. Then $f$ has a limiting distribution: $f(n)$, for $n$ uniform in $[1,x]$, converges in law as $x\to\infty$ to
\[
Z=\sum_pf(p^{G_p}),
\]
a series of independent variables that converges almost surely. Its characteristic function is
\[
\mathbb Ee^{itZ}=\prod_p\Bigl(1-\frac1p\Bigr)\Bigl(1+\sum_{k\ge1}\frac{e^{itf(p^k)}}{p^k}\Bigr).
\]
\end{theorem}

\begin{proof}
\emph{The model converges.} Let $X_p=f(p^{G_p})$ and apply the three-series theorem with $c=1$. Since $P(G_p\ge2)=p^{-2}$,
\[
\begin{gathered}P(|X_p|>1)\le\frac1{p^2}+\frac{\mathbf 1_{|f(p)|>1}}p,\\\Bigl|\mathbb E[X_p\mathbf 1_{|X_p|\le1}]-f^*(p)\frac{p-1}{p^2}\Bigr|\le\frac1{p^2},\qquad\mathbb E[X_p^2\mathbf 1_{|X_p|\le1}]\le\frac{f^*(p)^2}p+\frac1{p^2}.\end{gathered}
\]
The three hypotheses, with $\sum_p\frac1{p^2}<\infty$ and $|\frac{f^*(p)}{p^2}|\le\frac1{p^2}$, make the three series of Tome~10, Theorem~2.5, converge (the variance is at most the second moment). So $Z_y\to Z$ almost surely as $y\to\infty$, hence in law. Since $Z$ is the almost sure limit of independent sums, $\mathbb Ee^{itZ}=\lim_y\prod_{p\le y}\mathbb Ee^{itX_p}$ by dominated convergence, which is the product in the statement.

\emph{The small primes.} Let $y=x^{1/\log\log x}$, so $u=\log\log x\to\infty$ and $y\to\infty$. By Corollary~\ref{cor:kubiliusf} the laws of $f_y(n)$ and $Z_y$ are $o(1)$ apart in total variation. Since $Z_y\to Z$ in law, so does $f_y(n)$: for every bounded continuous $\phi$, $|\mathbb E_x\phi(f_y(n))-\mathbb E\phi(Z_y)|\le2\|\phi\|_\infty d_{TV}\to0$.

\emph{The large primes.} Let $R(n)=f(n)-f_y(n)=\sum_{p^k\parallel n,\,p>y}f(p^k)$. We show that $R(n)\to0$ in probability; Slutsky's lemma (Tome~3, Exercise~6.1) then gives $f(n)=f_y(n)+R(n)\to Z$ in law. Let $\tilde f$ be the strongly additive function with $\tilde f(p)=f^*(p)$ for $y<p\le x$ and $\tilde f(p)=0$ otherwise. Then $R(n)=\tilde f(n)$ unless $p^2\mid n$ for some $p>y$, or $p\mid n$ for some $p>y$ with $|f(p)|>1$. So
\[
P_x(R(n)\ne\tilde f(n))\le\sum_{p>y}\frac1{p^2}+\sum_{\substack{p>y\\|f(p)|>1}}\frac1p+O\Bigl(\frac{\pi(x)}x\Bigr)\longrightarrow0,
\]
as tails of convergent series (the last term accounts for the integer parts $\lfloor x/p\rfloor$ divided by $\lfloor x\rfloor$). By Tome~5, Theorem~9.1, applied to $\tilde f$, with $\tilde A=\sum_{y<p\le x}\frac{f^*(p)}p$ and $\tilde B^2=\sum_{y<p\le x}\frac{f^*(p)^2}p$,
\[
P_x\bigl(|\tilde f(n)-\tilde A|>\varepsilon\bigr)\le\frac{C_0\tilde B^2}{\varepsilon^2}\Bigl(1+O\Bigl(\frac1x\Bigr)\Bigr)\longrightarrow0,
\]
and $\tilde A\to0$, $\tilde B\to0$ as tails of the second and third series. So $P_x(|R(n)|>2\varepsilon)\to0$ for every $\varepsilon>0$.
\end{proof}

\begin{remark}[the converse]
Erd\H{o}s and Wintner also proved the converse: if $f$ has a limiting distribution, the three series converge. The proof needs a lower bound for characteristic functions of additive functions that we do not develop; see Tenenbaum's book, Chapter~III.4, or Elliott's \emph{Probabilistic Number Theory~I}, Chapter~5. Notice the asymmetry: the prime powers $p^k$ with $k\ge2$ play no role in the criterion, because $\sum_p\sum_{k\ge2}p^{-k}<\infty$.
\end{remark}

For $f(n)=\log\frac{\varphi(n)}n$ we have $f(p)=\log(1-\frac1p)\in[-\log2,0[$, so $f^*(p)=f(p)=-\frac1p+O(\frac1{p^2})$, and the three series converge (the first is empty). The same holds for $\log\frac{\sigma(n)}n$, with $f(p)=\log(1+\frac1p)$. So both have limiting distributions, and so do $\frac{\varphi(n)}n$ and $\frac{\sigma(n)}n$, by continuity of $\exp$. Figure~\ref{fig:phi} shows the first.

\begin{figure}[ht]
\centering\fig{ch01-phi}
\caption{The proportion of $n\le10^6$ with $\frac{\varphi(n)}n\le z$, against $z$. The limit law exists by Theorem~\ref{thm:erdoswintner} and is continuous by Corollary~\ref{cor:phicont}; the graph has no jumps, but its steep climbs at $z=\frac12$ (the even numbers not divisible by~$3$), $z=\frac13$ (multiples of $6$ prime to~$5$) and their neighbours show how singular it is. The mean is $\frac6{\pi^2}=0.6079$ (Exercise~\ref{ex:meanphi}).}
\label{fig:phi}
\end{figure}

\section{Continuity and purity of the limit law}

A sum of independent discrete variables is discrete when it has finitely many terms. With infinitely many it may be continuous: $\sum_k\varepsilon_k2^{-k}$, with independent fair $\varepsilon_k\in\{0,1\}$, is uniform on $[0,1]$. Paul L\'evy found in 1931 exactly when this happens.

\begin{think}
The limit law of $\frac{\varphi(n)}n$ is the law of $\prod_p(1-\frac1p)^{B_p}$, with independent $B_p$ equal to $1$ with probability about $\frac1p$. Does it have an atom, a value taken with positive probability? Compare with $\omega$-like functions where $f(p)=0$ for most primes.
\end{think}

The key is a lemma from the geometry of sums of signs.

\begin{lemma}[Littlewood--Offord, in Erd\H{o}s's form, 1945]\label{lem:lo}
Let $a_1,\dots,a_m$ be non-zero real numbers and $\varepsilon_1,\dots,\varepsilon_m$ independent with $P(\varepsilon_i=1)=P(\varepsilon_i=-1)=\frac12$. Then for every real $c$,
\[
P\Bigl(\sum_{i=1}^m\varepsilon_ia_i=c\Bigr)\le\binom m{\lfloor m/2\rfloor}2^{-m}\le\frac1{\sqrt m}.
\]
\end{lemma}
\begin{proof}
Replacing $(\varepsilon_i,a_i)$ by $(-\varepsilon_i,-a_i)$ does not change the law, so we may assume every $a_i>0$. A choice of signs is a subset $I=\{i:\varepsilon_i=1\}$ of $\{1,\dots,m\}$, and $\sum\varepsilon_ia_i=2\sum_{i\in I}a_i-\sum_ia_i$. If $I\subsetneq J$, the sum for $J$ exceeds the sum for $I$ by $2\sum_{J\setminus I}a_i>0$. So the sets $I$ giving the value $c$ form an \emph{antichain}: none contains another.

\emph{Sperner's bound.} An antichain $\mathcal F$ of subsets of an $m$-set has at most $\binom m{\lfloor m/2\rfloor}$ members. Indeed, consider the $m!$ maximal chains $\emptyset\subset\{\sigma(1)\}\subset\{\sigma(1),\sigma(2)\}\subset\cdots$, one for each permutation $\sigma$. A set $I$ of size $j$ lies on exactly $j!(m-j)!$ of them, and each chain meets the antichain at most once. Hence $\sum_{I\in\mathcal F}|I|!(m-|I|)!\le m!$, that is $\sum_{I\in\mathcal F}\binom m{|I|}^{-1}\le1$, and each $\binom m{|I|}\le\binom m{\lfloor m/2\rfloor}$.

\emph{The binomial bound.} For $m=2k$, $\bigl(\binom{2k}k4^{-k}\bigr)^2=\prod_{j=1}^k\bigl(\frac{2j-1}{2j}\bigr)^2\le\prod_{j=1}^k\frac{2j-1}{2j}\cdot\frac{2j}{2j+1}=\frac1{2k+1}$, so $\binom{2k}k4^{-k}\le\frac1{\sqrt{2k+1}}\le\frac1{\sqrt m}$. For $m=2k+1$, $\binom{2k+1}k2^{-2k-1}=\binom{2k}k4^{-k}\frac{2k+1}{2k+2}\le\frac1{\sqrt{2k+1}}=\frac1{\sqrt m}$.
\end{proof}

\begin{theorem}[L\'evy, 1931]\label{thm:levyatoms}
Let $X_1,X_2,\dots$ be independent discrete real random variables such that $S=\sum_jX_j$ converges almost surely, and let $q_j=\max_vP(X_j=v)$. Then the law of $S$ has an atom if and only if $\sum_j(1-q_j)<\infty$; otherwise it is continuous.
\end{theorem}

\begin{proof}
\emph{Convergent case.} Choose $v_j$ with $P(X_j=v_j)=q_j$. By the first Borel--Cantelli lemma, almost surely $X_j=v_j$ for all large $j$. Since $\sum X_j$ converges almost surely, $\sum_jv_j$ converges (its tail coincides with a tail of $\sum X_j$ on an event of positive probability). The events $E_N=\{X_j=v_j\ \forall j>N\}$ increase to an almost sure event, so $P(E_N)>0$ for some $N$. The finite sum $S_N=X_1+\dots+X_N$ is discrete, so $P(S_N=w)>0$ for some $w$. By independence,
\[
P\Bigl(S=w+\sum_{j>N}v_j\Bigr)\ge P(S_N=w)P(E_N)>0.
\]
\emph{Divergent case.} Let $S'$ be an independent copy of $S$, built from an independent copy $(X'_j)$. For every $a$,
\[
P(S=a)^2=P(S=a,S'=a)\le P(S-S'=0).
\]
Now $S-S'=\sum_jD_j$ with $D_j=X_j-X'_j$ independent and symmetric, and $P(D_j\ne0)=1-\sum_vP(X_j=v)^2\ge1-q_j$ since $\sum_vP(X_j=v)^2\le q_j\sum_vP(X_j=v)=q_j$. So $\sum_jP(D_j\ne0)=\infty$, and by the second Borel--Cantelli lemma infinitely many $D_j$ are non-zero, almost surely. Let $(\varepsilon_j)$ be independent fair signs, independent of everything else. By symmetry and independence, the sequence $(\varepsilon_jD_j)$ has the same law as $(D_j)$, so $P(\sum_jD_j=0)=P(\sum_j\varepsilon_jD_j=0)$.

Fix $m\ge1$ and $N$ so large that the event $F$ that at least $m$ of $D_1,\dots,D_N$ are non-zero has $P(F)\ge1-\frac1m$. Condition on all the $D_j$ and on the $\varepsilon_j$ with $D_j=0$ or $j>N$. On $F$, the sum $\sum_j\varepsilon_jD_j$ is then $c+\sum_{j\in J}\varepsilon_jD_j$ for a fixed $c$ and a set $J$ of at least $m$ indices with $D_j\ne0$, and the remaining signs are still independent and fair. Lemma~\ref{lem:lo} gives a conditional probability at most $\frac1{\sqrt m}$ that the sum is $0$. So $P(\sum D_j=0)\le\frac1{\sqrt m}+\frac1m$ for every $m$, hence $P(\sum D_j=0)=0$ and $S$ has no atom.
\end{proof}

\begin{corollary}\label{cor:ewcont}
Under the hypotheses of Theorem~\ref{thm:erdoswintner}, the limit law of $f$ is continuous if and only if $\sum_{f(p)\ne0}\frac1p=\infty$.
\end{corollary}
\begin{proof}
Let $X_p=f(p^{G_p})$. For $p\ge3$ every value other than $0$ is taken with probability at most $P(G_p\ge1)=\frac1p<1-\frac1p\le P(X_p=0)$, so $q_p=P(X_p=0)$ and $1-q_p=P(X_p\ne0)$. If $f(p)\ne0$ then $P(X_p\ne0)\ge P(G_p=1)=\frac{p-1}{p^2}\ge\frac1{2p}$; and always $P(X_p\ne0)\le\frac{\mathbf 1_{f(p)\ne0}}p+\frac1{p^2}$. So $\sum_p(1-q_p)$ and $\sum_{f(p)\ne0}\frac1p$ converge or diverge together, and Theorem~\ref{thm:levyatoms} applies.
\end{proof}

\begin{corollary}[Schoenberg 1928, Davenport 1933]\label{cor:phicont}
The functions $\frac{\varphi(n)}n$ and $\frac{\sigma(n)}n$ have continuous limiting distribution functions. In particular the set of abundant numbers, $\{n:\sigma(n)\ge2n\}$, has a natural density.
\end{corollary}
\begin{proof}
For both logarithms $f(p)\ne0$ for every $p$, and $\sum\frac1p=\infty$. Continuity of the limit law $F$ of $\log\frac{\sigma(n)}n$ at $\log2$ means that $P_x(\log\frac{\sigma(n)}n\ge\log2)\to1-F(\log2^-)=1-F(\log2)$.
\end{proof}

The argument shows more: the limit law of $\log\frac{\sigma(n)}n$ gives zero mass to every single point, so the density of $\{n:\sigma(n)\ge\alpha n\}$ exists for every $\alpha$ and depends continuously on $\alpha$. Figure~\ref{fig:sigma} shows the histogram up to $10^6$: $24.75\%$ of the integers are abundant there. The limit is $0.24761\ldots$, computed by Deléglise and then, to more digits, by Kobayashi and others, from the characteristic function of Theorem~\ref{thm:erdoswintner}.

\begin{figure}[ht]
\centering\fig{ch01-sigma}
\caption{Histogram of $\frac{\sigma(n)}n$ for $n\le10^6$ (density scale). The spike at $\frac{\sigma(n)}n=1^+$ comes from the primes and the integers with only large prime factors; abundant numbers are those right of the red line. The largest value below $10^6$ is $4.509$, at $n=720720$.}
\label{fig:sigma}
\end{figure}

Continuity is not the end of the story. A probability law on $\R$ splits uniquely as $\mu=\mu_d+\mu_{sc}+\mu_{ac}$, a discrete part, a continuous part carried by a Lebesgue-null set, and an absolutely continuous part (Tome~6, Theorem~10.5; the singular part then splits into its atoms and the rest). Jessen and Wintner proved in 1935 that laws of convergent sums of independent discrete variables are always of one pure type.

\begin{theorem}[Jessen and Wintner, 1935]\label{thm:jessenwintner}
Under the hypotheses of Theorem~\ref{thm:levyatoms}, the law $\mu$ of $S$ is either purely discrete, or purely singular continuous, or absolutely continuous.
\end{theorem}
\begin{proof}
Let $G$ be the additive subgroup of $\R$ generated by all the values $v$ with $P(X_j=v)>0$, for all $j$; it is countable. For a Borel set $B$, the set $B+G=\bigcup_{g\in G}(B+g)$ is Borel. Almost surely every partial sum $S_N$ lies in $G$, so $S\in B+G$ if and only if $\sum_{j>N}X_j\in B+G$. This holds for every $N$, so up to a null set $\{S\in B+G\}$ is a tail event, and $\mu(B+G)\in\{0,1\}$ by Kolmogorov's zero--one law.

If $\mu_d\ne0$, let $D$ be the countable set of atoms; then $\mu(D+G)\ge\mu(D)>0$, so $\mu(D+G)=1$, and $D+G$ is countable: $\mu$ is purely discrete. If $\mu_d=0$ and $\mu_{sc}\ne0$, let $N$ be a Lebesgue-null set with $\mu_{sc}(N)=\mu_{sc}(\R)$. Then $N+G$ is null, a countable union of null sets, so $\mu_{ac}(N+G)=0$; and $\mu(N+G)\ge\mu_{sc}(N)>0$ forces $\mu(N+G)=1$, hence $\mu_{ac}(\R)=\mu_{ac}(\R)-\mu_{ac}(N+G)\le\mu(\R\setminus(N+G))=0$. So $\mu=\mu_{sc}$. Otherwise $\mu=\mu_{ac}$.
\end{proof}

\begin{remark}[which pure type?]
Erd\H{o}s showed in 1939 that the limit laws of $\frac{\varphi(n)}n$ and $\frac{\sigma(n)}n$ are purely singular: their distribution functions are continuous but have derivative $0$ almost everywhere. The steep climbs in Figure~\ref{fig:phi} are the visible trace of this. Deciding the type of the law of $\sum_p f(p^{G_p})$ in general is hard, and is related to the theory of Bernoulli convolutions; see Elliott, \emph{Probabilistic Number Theory~I}, Chapter~7.
\end{remark}

\begin{beauty}
Two unrelated-looking facts, that one integer in four is abundant and that $\frac{\varphi(n)}n$ has a continuous distribution, are instances of one probabilistic theorem about sums of independent coins. The arithmetic enters only through two counts: a prime $p$ divides a proportion $\frac1p$ of the integers, and the divisibilities by different small primes are independent up to an error that the sieve controls. Everything else is probability, and it is the probability of Kolmogorov, L\'evy and Sperner, a combinatorialist's lemma about subsets.
\end{beauty}

\section{Exercises}

\begin{exercise}[$\star$]\label{ex:meanphi}
Show that $\sum_{n\le x}\frac{\varphi(n)}n=\frac6{\pi^2}x+O(\log x)$, using $\frac{\varphi(n)}n=\sum_{d\mid n}\frac{\mu(d)}d$. Check that $\frac6{\pi^2}$ is the mean of the limit law, $\prod_p\mathbb E[(1-\frac1p)^{\mathbf 1_{G_p\ge1}}]$.
\end{exercise}

\begin{exercise}[$\star$]\label{ex:omeganolimit}
Show that $\omega$ has no limiting distribution: for every $K$, $P_x(\omega(n)\le K)\to0$. Which of the three series diverges?
\end{exercise}

\begin{exercise}[$\star$]\label{ex:abundant6}
Show that every multiple of an abundant number is abundant, that $6k$ is abundant for $k\ge2$, and deduce that the density of abundant numbers is at least $\frac16+\frac1{20}-\frac1{60}=0.2$. (Use also $20$.)
\end{exercise}

\begin{exercise}[$\star\star$]\label{ex:phimoments}
Show that for every real $s>0$, $\frac1x\sum_{n\le x}(\frac{\varphi(n)}n)^s\to\prod_p\bigl(1-\frac1p+\frac1p(1-\frac1p)^s\bigr)$. Compute the limit for $s=1$ and $s=2$ as products, and evaluate them numerically.
\end{exercise}

\begin{exercise}[$\star\star$]\label{ex:purity}
Let $\varepsilon_k$ be independent with $P(\varepsilon_k=0)=P(\varepsilon_k=1)=\frac12$. Show that $S=\sum_{k\ge1}\varepsilon_k2^{-k}$ is uniform on $[0,1]$, and that $T=\sum_{k\ge1}2\varepsilon_k3^{-k}$ is continuous and carried by the Cantor set, hence purely singular. Check both against Theorems~\ref{thm:levyatoms} and~\ref{thm:jessenwintner}.
\end{exercise}

\begin{exercise}[$\star\star$]\label{ex:losharp}
Show that Lemma~\ref{lem:lo} is sharp up to a constant: for $a_1=\dots=a_m=1$, $\max_cP(\sum\varepsilon_ia_i=c)\sim\sqrt{\frac2{\pi m}}$. Show that if $a_1<\dots<a_m$ are distinct positive integers and $m\ge2$, then $P(\sum\varepsilon_ia_i=c)\le\frac12$, and compute this probability for $a=(1,2,3,4)$ and $c=0$.
\end{exercise}

\begin{exercise}[$\star\star$]\label{ex:kubiliusomega}
Use Corollary~\ref{cor:kubiliusf} to show that for $y=x^{1/\log\log x}$, the law of $\omega_y(n)=\#\{p\le y:p\mid n\}$ is within $o(1)$ in total variation of the law of a sum of independent Bernoulli variables with parameters $\frac1p$, $p\le y$. Deduce from the Chen--Stein bound of Tome~10 (Theorem~6.8) that $\omega_y(n)$ is within $O(\frac1{\log\log x})$ of a Poisson law of parameter $\sum_{p\le y}\frac1p$, up to the error of Theorem~\ref{thm:kubilius}.
\end{exercise}

\begin{exercise}[$\star\star\star$]\label{ex:phisupport}
Show that the support of the limit law of $\frac{\varphi(n)}n$ is the whole interval $[0,1]$: for every $0\le a<b\le1$, a positive proportion of integers have $a<\frac{\varphi(n)}n<b$. (Approximate a target $t\in\,]-\infty,0[$ by a finite sum of distinct terms $\log(1-\frac1p)$, using $\sum\frac1p=\infty$.)
\end{exercise}

\section{Hints}

\begin{hint}{ex:meanphi}
Swap the sums: $\sum_{n\le x}\sum_{d\mid n}\frac{\mu(d)}d=\sum_{d\le x}\frac{\mu(d)}d\lfloor\frac xd\rfloor$.
\end{hint}
\begin{hint}{ex:omeganolimit}
Tur\'an--Kubilius with Chebyshev's inequality.
\end{hint}
\begin{hint}{ex:abundant6}
$\frac{\sigma(n)}n=\sum_{d\mid n}\frac1d$ only grows when $n$ gets new divisors.
\end{hint}
\begin{hint}{ex:phimoments}
Bounded continuous functions of a variable that converges in law.
\end{hint}
\begin{hint}{ex:purity}
Binary and ternary expansions; for $T$, count the intervals of the $k$-th stage of the Cantor construction.
\end{hint}
\begin{hint}{ex:losharp}
Stirling's formula; for distinct integers, look at what happens when only the largest $a_i$ changes sign.
\end{hint}
\begin{hint}{ex:kubiliusomega}
$\omega_y(n)=\Lambda(v_y(n))$ with $\Lambda(\mathbf k)=\#\{p:k_p\ge1\}$, and $\mathbf 1_{G_p\ge1}$ is Bernoulli$(\frac1p)$.
\end{hint}
\begin{hint}{ex:phisupport}
A finite set of primes $p_1<\dots<p_r$, all larger than some $P$, and the integers $n$ with $p_i\mid n$ and no other prime factor in $[2,y]$ except those above $P'$.
\end{hint}

\section{Solutions}

\begin{solution}{ex:meanphi}
\idea Whenever an arithmetic function is a divisor sum $\sum_{d\mid n}g(d)$, its mean is $\sum_d\frac{g(d)}d$, by swapping the order of summation; the error is controlled by $\sum|g(d)|$.
\writeup Since $\varphi(n)=\sum_{d\mid n}\mu(d)\frac nd$, we have $\frac{\varphi(n)}n=\sum_{d\mid n}\frac{\mu(d)}d$, and
\[
\sum_{n\le x}\frac{\varphi(n)}n=\sum_{d\le x}\frac{\mu(d)}d\Bigl(\frac xd+O(1)\Bigr)=x\sum_{d\ge1}\frac{\mu(d)}{d^2}+O\Bigl(x\sum_{d>x}\frac1{d^2}\Bigr)+O\Bigl(\sum_{d\le x}\frac1d\Bigr)=\frac x{\zeta(2)}+O(\log x).
\]
In the model, $\frac{\varphi(n)}n$ becomes $\prod_p(1-\frac1p)^{\mathbf 1_{G_p\ge1}}$, a product of independent factors with mean $1-\frac1p+\frac1p(1-\frac1p)=1-\frac1{p^2}$. The product of the means is $\prod_p(1-p^{-2})=\frac1{\zeta(2)}=\frac6{\pi^2}$, and since $\frac{\varphi(n)}n$ is bounded the mean of the limit law is the limit of the means.
\end{solution}

\begin{solution}{ex:omeganolimit}
\idea If a function grows to infinity in probability, it cannot converge in law; Tur\'an--Kubilius shows that $\omega(n)$ is within $\sqrt{\log\log x}$ of $\log\log x$ for most $n$.
\writeup Here $f(p)=1$, so $f^*(p)=1$ and the second and third series are $\sum\frac1p=\infty$. With $A(x)=\sum_{p\le x}\frac1p=\log\log x+O(1)$ and $B(x)^2=A(x)$, Tome~5, Theorem~9.1, and Chebyshev give, for $x$ so large that $A(x)>2K$,
\[
P_x(\omega(n)\le K)\le P_x\Bigl(|\omega(n)-A|\ge\frac A2\Bigr)\le\frac{4C_0B^2}{A^2}(1+o(1))=\frac{4C_0}{A}(1+o(1))\to0.
\]
\end{solution}

\begin{solution}{ex:abundant6}
\idea Abundance only depends on the set of divisors, and multiplying $n$ adds divisors.
\writeup $\frac{\sigma(n)}n=\sum_{d\mid n}\frac1d$. If $n\mid N$, then $\frac{\sigma(N)}N=\sum_{e\mid N}\frac1e\ge\sum_{d\mid n}\frac1d=\frac{\sigma(n)}n$, since the divisors of $n$ are divisors of $N$. So multiples of abundant numbers are abundant. For $k\ge2$, $6k$ has the divisors $6k,3k,2k,k$ and $1$, so $\sigma(6k)\ge(6+3+2+1)k+1>12k$. The number $20$ has $\sigma(20)=42>40$. The multiples of $6$ or $20$ have density $\frac16+\frac1{20}-\frac1{\operatorname{lcm}(6,20)}=\frac16+\frac1{20}-\frac1{60}=\frac15$, and all but $6$ itself (which is perfect, $\sigma(6)=12$, hence abundant in our convention $\sigma(n)\ge2n$) are abundant. By Corollary~\ref{cor:phicont} the density exists, so it is at least $0.2$.
\end{solution}

\begin{solution}{ex:phimoments}
\idea Convergence in law means convergence of $\mathbb E\phi$ for bounded continuous $\phi$; $(\frac{\varphi(n)}n)^s$ is such a function of $\log\frac{\varphi(n)}n$, and the limit is computed in the independent model.
\writeup Let $\phi(t)=e^{st}$ for $t\le0$; it is bounded and continuous on $]-\infty,0]$, where $\log\frac{\varphi(n)}n$ lives. By Theorem~\ref{thm:erdoswintner}, $\mathbb E_x\phi(f(n))\to\mathbb E\phi(Z)$, and $e^{sZ}=\prod_p(1-\frac1p)^{s\mathbf 1_{G_p\ge1}}$ (strong additivity: only whether $G_p\ge1$ matters). The partial products converge almost surely and are bounded by $1$, so by dominated convergence and independence $\mathbb Ee^{sZ}=\prod_p(1-\frac1p+\frac1p(1-\frac1p)^s)$. For $s=1$ this is $\prod_p(1-p^{-2})=0.6079$. For $s=2$, $1-\frac1p+\frac1p(1-\frac2p+\frac1{p^2})=1-\frac2{p^2}+\frac1{p^3}$, and $\prod_p(1-\frac2{p^2}+\frac1{p^3})=0.4282$ (the notebook evaluates the product with an Euler-product acceleration).
\end{solution}

\begin{solution}{ex:purity}
\idea Both are random digit expansions; the base and the allowed digits decide whether every interval gets its fair share.
\writeup For $S$, the binary digits of a uniform variable on $[0,1]$ are independent fair bits (Tome~3), and conversely the law of $S$ gives each dyadic interval $[\frac j{2^k},\frac{j+1}{2^k}]$ mass $2^{-k}$, which determines the uniform law. For $T$, the digits in base $3$ are $0$ or $2$, so $T$ lies in the Cantor set $C$, which at stage $k$ is covered by $2^k$ intervals of length $3^{-k}$: its Lebesgue measure is $0$. Each term $2\varepsilon_k3^{-k}$ has $q_k=\frac12$ and $\sum(1-q_k)=\infty$, so $T$ has no atom by Theorem~\ref{thm:levyatoms}; a law without atoms carried by a null set is purely singular. Both are pure, as Theorem~\ref{thm:jessenwintner} predicts; the two series differ only in the ratio $2$ versus $3$, which shows that the type cannot be read off from L\'evy's criterion alone.
\end{solution}

\begin{solution}{ex:losharp}
\idea For equal coefficients the sum is a centred binomial, whose largest probability is the central binomial coefficient. For distinct coefficients, flipping one sign moves the sum, so the sign vectors can be paired off.
\writeup With all $a_i=1$, $\sum\varepsilon_i=2|I|-m$, and the most likely value has probability $\binom m{\lfloor m/2\rfloor}2^{-m}$, which is $\sim\sqrt{\frac2{\pi m}}$ by Stirling's formula. For distinct positive $a_i$, pair each sign vector with the one in which only $\varepsilon_m$ is flipped. The two sums differ by $2a_m\ne0$, so at most one of each pair hits $c$, and $P\le\frac12$ (this uses only $a_m\ne0$). For $a=(1,2,3,4)$ and $c=0$, the set $I$ of plus signs must have $\sum_{I}a_i=5$, so $I=\{1,4\}$ or $I=\{2,3\}$, and $P=\frac2{16}=\frac18$. In fact for distinct coefficients the truth is of order $m^{-3/2}$, a theorem of S\'ark\"ozy and Szemer\'edi (1965).
\end{solution}

\begin{solution}{ex:kubiliusomega}
\idea $\omega_y(n)$ is a function of the exponent vector, so it inherits the total variation bound; in the model it is a sum of independent indicators, for which Tome~10 has a Poisson approximation.
\writeup $\omega_y(n)=\Lambda(v_y(n))$ with $\Lambda(\mathbf k)=\#\{p\le y:k_p\ge1\}$, and $\Lambda(G^{(y)})=\sum_{p\le y}\mathbf 1_{G_p\ge1}$, a sum of independent Bernoulli variables $B_p$ with $P(B_p=1)=P(G_p\ge1)=\frac1p$. By Theorem~\ref{thm:kubilius} with $u=\log\log x$ the two laws differ by $O(e^{-u/4}+x^{-1/4}\log x)=O((\log x)^{-1/4})$ in total variation. The Chen--Stein bound for a sum of independent Bernoulli variables gives $d_{TV}(\sum B_p,\operatorname{Poisson}(\lambda))\le\frac{1-e^{-\lambda}}\lambda\sum_p\frac1{p^2}$ with $\lambda=\sum_{p\le y}\frac1p=\log\log x+O(\log\log\log x)$; this is $O(\frac1{\log\log x})$. The triangle inequality gives the claim. Chapter~3 shows that the full $\omega(n)$, which includes the large primes, is \emph{not} that close to a Poisson law: the large primes add a non-Poisson correction, the factor $\frac1{\Gamma(z)}$ of the Selberg--Delange theorem.
\end{solution}

\begin{solution}{ex:phisupport}
\idea To make $\frac{\varphi(n)}n$ close to a target, choose the set of primes dividing $n$; the divergence of $\sum\frac1p$ lets a greedy choice among large primes reach any target with any precision, and the Kubilius model shows that the integers with a prescribed set of small prime divisors have positive density.
\writeup Let $0<a<b\le1$ (the case $a=0$ follows by shrinking the interval), and fix $t$ with $\log a<t<\log b$, $t<0$. Let $\eta=\frac12\min(t-\log a,\log b-t)$ and choose $P$ with $|\log(1-\frac1p)|<\eta$ for $p>P$. Since $\sum_{p>P}|\log(1-\frac1p)|=\infty$ and each term is less than $\eta$, adding the terms one by one in increasing order of $p$ and stopping just before the total passes $|t|$ gives primes $P<p_1<\dots<p_r$ with $s=\sum_i\log(1-\frac1{p_i})\in\,]t-\eta,t]$. Next choose $Q>p_r$ so large that $\frac2\eta\sum_{p>Q}\frac1{p^2}<\frac12$. In the model, let $T=\sum_{p>Q}\log(1-\frac1p)\mathbf 1_{G_p\ge1}\le0$; since $|\log(1-\frac1p)|\le\frac2p$, Markov's inequality gives $P(|T|>\eta)\le\frac1\eta\sum_{p>Q}\frac2p\cdot\frac1p<\frac12$. The event $E$: ``$G_{p_i}\ge1$ for $i\le r$, $G_p=0$ for the other $p\le Q$, and $|T|\le\eta$'' has probability at least $\frac12\prod_i\frac1{p_i}\prod_{p\le Q}(1-\frac1p)>0$, by independence. On $E$, $\log Z=s+T\in[s-\eta,s]\subset\,]t-2\eta,t]\subseteq\,]\log a,\log b[$. So the limit law gives positive mass to $]a,b[$. By Theorem~\ref{thm:erdoswintner} and Corollary~\ref{cor:phicont} (the endpoints carry no mass), the proportion of $n\le x$ with $a<\frac{\varphi(n)}n<b$ tends to that positive mass.
\end{solution}

\section*{Sources and further reading}

\begin{itemize}
\item \emph{History.} I.~J.~Schoenberg, ``Über die asymptotische Verteilung reeller Zahlen mod~1'', \emph{Math.\ Z.}\ 28 (1928), 171--199, for $\frac{\varphi(n)}n$. H.~Davenport, ``Über numeri abundantes'', \emph{Sitzungsber.\ Preuss.\ Akad.\ Wiss.}\ (1933), 830--837. F.~Behrend, ``Über numeri abundantes II'', same \emph{Sitzungsber.}\ (1933), 280--293, for the bounds $0.241$ and $0.314$. M.~Deléglise, ``Bounds for the density of abundant integers'', \emph{Experiment.\ Math.}\ 7 (1998), 137--143; G.~McNew and S.~Setty, ``On the densities of covering numbers and abundant numbers'', arXiv:2507.23041 (2025), for the value $0.24761\ldots$
\item \emph{Erd\H{o}s--Wintner.} P.~Erd\H{o}s and A.~Wintner, ``Additive arithmetical functions and statistical independence'', \emph{Amer.\ J.\ Math.}\ 61 (1939), 713--721. P.~Erd\H{o}s, ``On the smoothness of the asymptotic distribution of additive arithmetical functions'', \emph{Amer.\ J.\ Math.}\ 61 (1939), 722--725.
\item \emph{Kubilius.} J.~Kubilius, \emph{Probabilistic Methods in the Theory of Numbers}, Translations of Mathematical Monographs~11, American Mathematical Society, 1964. Biographical facts from E.~Manstavičius, ``Jonas Kubilius (1921--2011)'', \emph{Acta Arith.}\ 157 (2013), 11--36, and the entry ``Jonas Kubilius'' of the \emph{Visuotinė lietuvių enciklopedija} (rector 1958--1990). P.~D.~T.~A.~Elliott, ``Jonas Kubilius and probabilistic number theory: some personal reflections'', \emph{Lith.\ Math.\ J.}\ 55 (2015), 2--24.
\item \emph{Probability.} P.~L\'evy, ``Sur les séries dont les termes sont des variables éventuelles indépendantes'', \emph{Studia Math.}\ 3 (1931), 119--155. B.~Jessen and A.~Wintner, ``Distribution functions and the Riemann zeta function'', \emph{Trans.\ Amer.\ Math.\ Soc.}\ 38 (1935), 48--88. P.~Erd\H{o}s, ``On a lemma of Littlewood and Offord'', \emph{Bull.\ Amer.\ Math.\ Soc.}\ 51 (1945), 898--902.
\item \emph{Textbooks.} G.~Tenenbaum, \emph{Introduction to Analytic and Probabilistic Number Theory}, 3rd ed., AMS, 2015, Chapters~III.2--III.4. P.~D.~T.~A.~Elliott, \emph{Probabilistic Number Theory~I}, Springer, 1979. E.~Kowalski, \emph{An Introduction to Probabilistic Number Theory}, Cambridge University Press, 2021, Chapter~2, whose presentation of the Kubilius model we follow.
\end{itemize}
